%% FullCourse_10Lecs -- Lecture 7, Topic 7.1b: Word Embeddings
%% Source: split from Lec07_T1_Tokenization_Embeddings.tex section 3
%% Pass B: layer in refinements from
%%         ../../MiniCourse_8hr/Slides/Module3_LLM_Text/M3_T1_Tokenization_Embeddings.tex
%%         (Analogies and bias, Macro application gallery)

%% Shared preamble for the FullCourse_10Lecs topic decks.
%%
%% Each topic .tex \input's this file, then defines its own \title, \date
%% override (if any), \graphicspath, and \begin{document}.
%%
%% Reconstructed verbatim from the inlined copy preserved in
%% Lec10_Case_Studies/Lec10_T8_Case6_DeepRamsey.tex (lines 23-190), which is
%% explicitly annotated there as "from ../shared_preamble.tex, copied verbatim".
%% Base preamble matches MiniCourse_8hr/Slides/shared_preamble.tex; this version
%% additionally provides \sectiondivider, the conceptbox/cautionbox/examplebox/
%% takeawaybox tcolorboxes, and \compacttables used across the split decks.

\documentclass[10pt,english,aspectratio=169]{beamer}

\usepackage{ae,aecompl}
\usepackage[T1]{fontenc}
\usepackage{textcomp}
\usepackage[utf8]{inputenc}
\usepackage{babel}
\usepackage{datetime2}

\setcounter{secnumdepth}{3}
\setcounter{tocdepth}{3}

\usepackage{amsmath, amssymb, amsthm, graphicx}
\usepackage{booktabs, multirow}
\usepackage{tikz}
\usepackage{pgfplots}
\pgfplotsset{compat=1.16}
\usepackage[most]{tcolorbox}
\tcbuselibrary{skins, breakable, theorems}
\usepackage{listings}
\usepackage{xcolor}

\lstset{
  basicstyle=\ttfamily\small,
  keywordstyle=\color{blue!80!black}\bfseries,
  commentstyle=\color{green!50!black}\itshape,
  stringstyle=\color{orange!80!black},
  backgroundcolor=\color{gray!8},
  frame=single,
  rulecolor=\color{gray!40},
  breaklines=true,
  showstringspaces=false,
  numbers=none,
}

\lstdefinestyle{pythonstyle}{
    language=Python,
    basicstyle=\ttfamily\footnotesize,
    keywordstyle=\color{blue!80!black}\bfseries,
    commentstyle=\color{green!50!black}\itshape,
    stringstyle=\color{orange!80!black},
    frame=single,
    breaklines=true,
    showstringspaces=false,
    numbers=left,
    numbersep=5pt,
    numberstyle=\tiny\color{gray},
    backgroundcolor=\color{gray!8},
    rulecolor=\color{gray!40}
}

\ifx\hypersetup\undefined
  \AtBeginDocument{%
    \hypersetup{bookmarks=false,breaklinks=true,pdfborder={0 0 0},colorlinks=true,
      linkcolor=DarkText,urlcolor=RedTitle}%
  }
\else
  \hypersetup{bookmarks=false,breaklinks=true,pdfborder={0 0 0},colorlinks=true,
    linkcolor=DarkText,urlcolor=RedTitle}%
\fi

\makeatletter
\newcommand\makebeamertitle{%
  \begin{frame}[plain]
    \vspace*{1.0cm}
    \centering
    {\usebeamerfont{title}\usebeamercolor[fg]{title}\inserttitle\par}
    \vspace{1.0cm}
    {\usebeamerfont{author}\usebeamercolor[fg]{author}\insertauthor\par}
    \vspace{0.2cm}
    {\usebeamerfont{date}\usebeamercolor[fg]{date}\insertdate\par}
  \end{frame}}%
\AtBeginDocument{%
  \let\origtableofcontents=\tableofcontents
  \def\tableofcontents{\@ifnextchar[{\origtableofcontents}{\gobbletableofcontents}}
  \def\gobbletableofcontents#1{\origtableofcontents}%
}

\usetheme{metropolis}
\usefonttheme{professionalfonts}

\definecolor{LightGray}{RGB}{230,230,230}
\definecolor{RedTitle}{RGB}{0,0,200}
\definecolor{VeryLightGray}{RGB}{245,245,245}
\definecolor{DarkText}{RGB}{0,0,0}

\setbeamercolor{background canvas}{bg=white}
\setbeamercolor{title}{fg=RedTitle}
\setbeamercolor{author}{fg=DarkText}
\setbeamercolor{date}{fg=DarkText}
\setbeamercolor{frametitle}{bg=VeryLightGray,fg=DarkText}
\setbeamercolor{normal text}{fg=DarkText}
\setbeamerfont{title}{size=\large}

\setbeamersize{text margin left=5mm, text margin right=5mm}
\addtobeamertemplate{frametitle}{}{\vspace{-0.5em}}
\newcommand{\reducevspace}{\vspace{-0.5cm}}

% Shared section divider and box styles for split topic decks.
\newcommand{\sectiondivider}[2]{%
\begin{frame}[plain,noframenumbering]
\begin{center}
\vfill
{\Large\textbf{#1}\par}
\vspace{0.35cm}
{\normalsize #2\par}
\vfill
\end{center}
\end{frame}
}

\newtcolorbox{conceptbox}[1][]{
  colback=blue!3,
  colframe=RedTitle!55,
  boxrule=0.35mm,
  arc=1mm,
  left=2mm,
  right=2mm,
  top=1mm,
  bottom=1mm,
  #1
}
\newtcolorbox{cautionbox}[1][]{
  colback=red!3,
  colframe=red!50!black,
  boxrule=0.35mm,
  arc=1mm,
  left=2mm,
  right=2mm,
  top=1mm,
  bottom=1mm,
  #1
}
\newtcolorbox{examplebox}[1][]{
  colback=green!3,
  colframe=green!45!black,
  boxrule=0.35mm,
  arc=1mm,
  left=2mm,
  right=2mm,
  top=1mm,
  bottom=1mm,
  #1
}
\newtcolorbox{takeawaybox}[1][]{
  colback=VeryLightGray,
  colframe=RedTitle!55,
  boxrule=0.35mm,
  arc=1mm,
  left=2mm,
  right=2mm,
  top=1mm,
  bottom=1mm,
  #1
}
\newcommand{\compacttables}{\renewcommand{\arraystretch}{1.15}}

\setbeamertemplate{navigation symbols}{}
\setbeamertemplate{footline}{%
  \hfill%
  \usebeamercolor[fg]{page number in head/foot}%
  \usebeamerfont{page number in head/foot}%
  \insertframenumber\,/\,\inserttotalframenumber\kern1em\vskip2pt%
}

\makeatother

\author{Zhigang Feng}
% "date with time": compile timestamp so each build is identifiable.
\date{\today\ \DTMcurrenttime}


\graphicspath{{./pic/}{../../MiniCourse_8hr/Slides/Module3_LLM_Text/pic/}{../}}

\title[AI for Econ Research]{AI for Economic Research: Dynamic Models, Language, and Agents\\
Lecture 7.1b: Word Embeddings}

\begin{document}

\makebeamertitle

%=============================================================================
\begin{frame}{Topic 7.1b Agenda}
\begin{enumerate}
  \item \textbf{Word Embeddings} --- TF-IDF $\to$ Word2Vec $\to$ GloVe
  \medskip
  \item \textbf{Geometry of meaning} --- Arrow--Debreu analogy, analogies and bias
  \medskip
  \item \textbf{From static to contextual} --- limits of static embeddings; macro applications
\end{enumerate}

\bigskip
\textit{Builds on Topic 7.1a (tokenization). Bridge to Topic 7.2 (self-attention and the Transformer).}

\medskip
\footnotesize\textit{Arc: T1a Why$+$Tokens $\to$ T1b Embeddings $\to$ T2a Attention $\to$ T2b Transformer $\to$ T3a Training $\to$ T3b Limits $\to$ T4 Applications.}
\end{frame}

%=============================================================================
\section{Word Embeddings}

\begin{frame}{Why Embed Tokens into Vectors?}
\begin{itemize}
\item Neural networks operate on real-valued vectors, not discrete symbols.

\medskip
\item An \textbf{embedding} is a learned function
\[
   \mathcal{E}\colon \{1, \ldots, M\} \;\longrightarrow\; \mathbb{R}^{d},
\]
mapping each token ID to a dense vector of dimension $d$ (typically $d=128$ to $4096$).

\medskip
\item \textbf{Think of it as a basis problem from linear algebra.} We want to find a coordinate system --- a set of $d$ orthogonal ``semantic axes'' (e.g.\ monetary-policy stance, positive valence, corporate context) --- such that every token maps to a point in $\mathbb{R}^{d}$. Once tokens are coordinates, we can do arithmetic: measure distances, compute cosine similarities, add and subtract meaning. We will later connect this to Arrow--Debreu securities in financial markets. The challenge: we do not know the right axes in advance. They must be \emph{learned} from data.

\medskip
\item \textbf{What we want from embeddings:}
\begin{itemize}
    \item Semantically similar tokens should have nearby vectors (close coordinates in the learned basis)
    \item Algebraic operations (sums, differences) should reflect semantic relations
    \item The mapping should be \emph{learned} from data, not hand-designed
\end{itemize}

\medskip
\item \textbf{Bonus:} we will also need a \emph{positional} signal so the model knows token order---more on this when we meet the Transformer.
\end{itemize}
\end{frame}

\begin{frame}{Three Approaches to Word Representation}
\textbf{The problem from the previous slide:} find a mapping from tokens to vectors that captures semantic relationships. Three generations of methods have tried to solve this, each fixing a specific failure of its predecessor. Understanding this progression explains \emph{why} the Transformer was needed.
\medskip
\begin{center}
\renewcommand{\arraystretch}{1.3}
\begin{tabular}{l l l l}
\toprule
\textbf{Method} & \textbf{Type} & \textbf{Idea} & \textbf{Year} \\
\midrule
TF-IDF & Sparse, frequency & Re-weight word counts & 1970s \\
Word2Vec & Dense, predictive & Predict neighbors in window & 2013 \\
GloVe & Dense, statistical & Factor co-occurrence matrix & 2014 \\
Transformer embeddings & Dense, contextual & Vary by surrounding tokens & 2017+ \\
\bottomrule
\end{tabular}
\end{center}

\smallskip
{\footnotesize\textit{Transformer row previewed for completeness; full treatment in T2a/T2b. This topic covers only the first three rows.}}

\medskip
\begin{itemize}
\item TF-IDF: a bag-of-words feature. Transparent and easy to interpret.
\item Word2Vec / GloVe: produce a single vector per word \emph{type} (\emph{static} embeddings).
\item Transformer embeddings: produce a different vector for each \emph{occurrence} of a word, depending on its surrounding context (\emph{contextual} embeddings).
\end{itemize}

\medskip
We'll meet each in turn, then see why static embeddings break down and motivate the Transformer.
\end{frame}

\begin{frame}{Term Frequency (TF)}
\begin{itemize}
\item \textbf{Definition:} share of document $d$ consisting of word $w$:
\[
   \text{TF}(w, d) \;=\; \frac{\text{count}(w, d)}{\text{total words in } d}.
\]

\medskip
\item \textbf{Toy example.} FOMC-like sentence $d$: \emph{``the Fed raised rates because inflation persisted''} (7 words). Every word appears once, so TF is $1/7 \approx 0.14$ for each:
\begin{center}
\begin{tabular}{l c c}
\toprule
word & count in $d$ & TF$(w,d)$ \\
\midrule
\emph{the} & 1 & $0.14$ \\
\emph{Fed} & 1 & $0.14$ \\
\emph{raised} & 1 & $0.14$ \\
\emph{rates} & 1 & $0.14$ \\
\emph{inflation} & 1 & $0.14$ \\
\bottomrule
\end{tabular}
\end{center}

\medskip
\item \textbf{Limitation:} TF rewards \emph{frequent} words. Across a corpus, function words (\emph{the, a, is}) dominate even though they carry almost no semantic signal. We need a counterweight.
\end{itemize}
\end{frame}

\begin{frame}{Inverse Document Frequency (IDF)}
\begin{itemize}
\item \textbf{Definition:} how rare word $w$ is across a collection of $N$ documents:
\[
   \text{IDF}(w) \;=\; \log\!\left(\frac{N}{\#\{\text{docs containing } w\}}\right).
\]

\medskip
\item \textbf{Toy example.} Corpus of $N = 100$ FOMC speeches:
\begin{center}
\begin{tabular}{l c c}
\toprule
word & docs containing it & IDF$(w)$ \\
\midrule
\emph{the}         & 100 & $\log(100/100) = 0$ \\
\emph{inflation}   & 60  & $\log(100/60) \approx 0.51$ \\
\emph{yield curve} & 5   & $\log(100/5) \approx 3.00$ \\
\bottomrule
\end{tabular}
\end{center}

\medskip
\item \textbf{Intuition:} rare $\Rightarrow$ informative $\Rightarrow$ up-weighted. \emph{``the''} gets $0$ automatically --- no stop-word list needed. \emph{``yield curve''} gets a large multiplier because it distinguishes a few speeches from the rest.
\end{itemize}
\end{frame}

\begin{frame}{Putting It Together: TF-IDF}
\begin{itemize}
\item \textbf{Product formula:} each (word, document) pair gets a weighted score
\[
   \text{TF-IDF}(w, d) \;=\; \text{TF}(w, d)\,\cdot\,\text{IDF}(w).
\]
Stacking across all words in vocabulary $V$ $\Rightarrow$ each document is a sparse vector in $\mathbb{R}^{|V|}$.

\medskip
\item \textbf{Toy 3-speech matrix} (illustrative scores):
\begin{center}
\begin{tabular}{l c c c}
\toprule
 & \emph{the} & \emph{inflation} & \emph{yield curve} \\
\midrule
speech A (hawkish)   & 0 & 0.08 & 0.00 \\
speech B (dovish)    & 0 & 0.05 & 0.00 \\
speech C (technical) & 0 & 0.02 & 0.21 \\
\bottomrule
\end{tabular}
\end{center}

\medskip
\item \textbf{Properties:}
\begin{itemize}
    \item Down-weights common words automatically (no stop-word list)
    \item Highlights distinctive words for each document
    \item Lives in $\mathbb{R}^{|V|}$ --- very high-dimensional, very sparse
    \item \emph{No} semantic similarity: \emph{``central bank''} and \emph{``Federal Reserve''} are orthogonal $\Rightarrow$ motivation for dense embeddings (next slide).
\end{itemize}
\end{itemize}
\end{frame}

\begin{frame}{Word2Vec: Dense Embeddings from Context}
\begin{itemize}
\item \textbf{Distributional hypothesis} (Firth, 1957): \emph{``You shall know a word by the company it keeps.''}

\medskip
\item \textbf{Idea} (Mikolov et al., 2013): for each word in a corpus, look at a window of $\pm L$ neighboring words. Learn a vector for the center word that \emph{predicts} its neighbors (\emph{Skip-gram}) or learn one that \emph{is predicted by} its neighbors (\emph{CBOW}).

\medskip
\item \textbf{Training output:} a single static vector $\mathbf{u}_w \in \mathbb{R}^{d}$ for each word $w$ in the vocabulary, with $d$ typically $100$--$300$.

\medskip
\item \textbf{Geometric magic:} after training, simple vector arithmetic captures analogies, e.g.
\[
   \mathbf{u}_{\text{king}} - \mathbf{u}_{\text{man}} + \mathbf{u}_{\text{woman}} \;\approx\; \mathbf{u}_{\text{queen}}.
\]
Words living on the same semantic axis differ by approximately the same vector.
\end{itemize}
\end{frame}

\begin{frame}{Skip-Gram in One Slide: The Math}
\begin{itemize}
\item Each word $w$ has \emph{two} learned vectors: an input vector $\mathbf{u}_{w}$ and an output vector $\mathbf{v}_{w}$ (often merged after training).

\medskip
\item For a center word $w_t$ and a context word $w_c$ in its window, the model predicts
\[
   P\bigl(w_c \mid w_t\bigr) \;=\; \frac{\exp\bigl(\mathbf{u}_{w_t}^{\top} \mathbf{v}_{w_c}\bigr)}{\sum_{j=1}^{M}\exp\bigl(\mathbf{u}_{w_t}^{\top} \mathbf{v}_{w_j}\bigr)}.
\]

\medskip
\item Training objective (negative log-likelihood over all (center, context) pairs):
\[
   \mathcal{L} \;=\; -\sum_{t}\sum_{c \in \text{window}(t)} \log P\bigl(w_c \mid w_t\bigr).
\]

\medskip
\item Optimized with stochastic gradient descent (in practice with \emph{negative sampling} to avoid the full vocabulary softmax).

\medskip
\item \textbf{Result:} words with similar context distributions end up with similar $\mathbf{u}$ vectors.
\end{itemize}
\end{frame}

\begin{frame}{Skip-Gram: Walkthrough on ``central banks target inflation''}
\begin{itemize}
\item \textbf{Sentence:} \emph{``central banks target inflation''}, window radius $L = 2$, center word $=$ \emph{target}.

\medskip
\item Training pairs formed from this center word:
\[
(\textit{target}, \textit{central}), \quad (\textit{target}, \textit{banks}), \quad (\textit{target}, \textit{inflation}).
\]

\medskip
\item For each pair, the score is the dot product $\mathbf{u}_{\text{target}}^{\top} \mathbf{v}_{\text{context}}$. Gradient descent \emph{increases} this score for observed pairs and \emph{decreases} it for random (negative) pairs.

\medskip
\item \textbf{Across the whole corpus:} \emph{target} appears near \emph{rates}, \emph{inflation}, \emph{policy}, \emph{Fed}. So do \emph{raise} and \emph{cut}. Consequently
\[
\mathbf{u}_{\text{target}} \;\text{ends up close to}\; \mathbf{u}_{\text{raise}},\, \mathbf{u}_{\text{cut}} \;\text{and far from}\; \mathbf{u}_{\text{apple}}.
\]

\medskip
\item \textbf{Takeaway:} analogies emerge because consistent co-occurrences place related words at similar positions in $\mathbb{R}^{d}$.
\end{itemize}
\end{frame}

\begin{frame}{GloVe: Global Co-occurrence Statistics}
\begin{itemize}
\item Pennington et al.\ (2014): instead of sliding-window prediction, factor a \emph{global} matrix.

\medskip
\item Let $X_{ij}$ count how often word $j$ appears near word $i$ across the entire corpus. GloVe fits embeddings to satisfy
\[
   \mathbf{u}_{w_i}^{\top}\,\mathbf{v}_{w_j} + b_{i} + b_{j} \;\approx\; \log X_{ij}.
\]

\medskip
\item \textbf{Toy counts on a central-bank corpus} $X_{ij}$ (illustrative):
\begin{center}
\small
\renewcommand{\arraystretch}{1.1}
\begin{tabular}{l c c c}
\toprule
 & \emph{rates} & \emph{inflation} & \emph{apple} \\
\midrule
Fed      & 820 & 410 & 3 \\
ECB      & 540 & 390 & 2 \\
hawkish  & 210 & 160 & 0 \\
\bottomrule
\end{tabular}
\end{center}
$\log X_{\text{Fed},\text{rates}} \approx 6.7$ (large target dot-product) vs.\ $\log X_{\text{Fed},\text{apple}} \approx 1.1$ (small). Fitting the equation therefore forces $\mathbf{u}_{\text{Fed}}$ close to $\mathbf{v}_{\text{rates}}, \mathbf{v}_{\text{inflation}}$ and orthogonal to $\mathbf{v}_{\text{apple}}$ --- monetary-policy structure falls out of corpus-wide statistics.

\medskip
\item \textbf{Contrast with Word2Vec:} predictive / local / streaming (W2V) vs.\ statistical / global / one-shot factorization (GloVe). Both produce \emph{static} dense vectors with similar geometry. Pick W2V for massive streaming corpora, GloVe when a full co-occurrence matrix is cheap to build.
\end{itemize}
\end{frame}

%=============================================================================
\section{Geometry of Meaning}

\begin{frame}{A 3D Toy of Word Embeddings}
\begin{center}
\includegraphics[width=10cm,height=5.8cm]{word_examples.png}
\par\end{center}
\small
The three axes loosely correspond to ``Animal-ness'', ``Royalty'', and ``Fruit-ness''. Words like \emph{cat}, \emph{king}, and \emph{apple} cluster along the dimension that fits them best. Real LLM embeddings live in hundreds of dimensions, but the same intuition holds.
\end{frame}

\begin{frame}{Analogy: Embeddings as Arrow--Debreu Securities}
\begin{itemize}
\item A $d$-dimensional embedding is the linguistic cousin of an Arrow--Debreu payoff vector: the $d$ coordinates are payoffs across $d$ latent semantic ``states.''

\medskip
\item \textbf{Dictionary, term by term:}
\begin{itemize}
    \item \emph{$d$-dim embedding} $\equiv$ payoff vector across $d$ latent semantic states (royalty, monetary-policy stance, animal-ness, $\ldots$).
    \item \emph{Inner product} $\mathbf{u}^{\top}\mathbf{v}$ $\equiv$ covariance of state-payoffs; \emph{orthogonality} $\Leftrightarrow$ uncorrelated semantic content.
    \item \emph{Cosine similarity} $\cos(\mathbf{u}, \mathbf{v}) = \tfrac{\mathbf{u}^{\top}\mathbf{v}}{\|\mathbf{u}\|\,\|\mathbf{v}\|}$ $\equiv$ correlation under unit-norm payoffs --- length normalization removes the influence of an asset's overall scale.
    \item \emph{Arrow--Debreu prices over states} $\leftrightarrow$ \emph{softmax attention weights over tokens} (preview of T2a) --- both are non-negative weights summing to 1 over a state space, and both price/aggregate state-contingent payoffs.
\end{itemize}

\medskip
\item \textbf{The pedagogical pay-off.} This is why economists already have an intuition for everything in T2a. Self-attention is just a learned, content-addressed state-pricing kernel.
\end{itemize}
\end{frame}

\begin{frame}{Analogies and the Bias Problem}
\begin{itemize}
\item \textbf{Geometric magic} (Mikolov et al., 2013): simple vector arithmetic captures analogies.
\[
   \mathbf{u}_{\text{king}} - \mathbf{u}_{\text{man}} + \mathbf{u}_{\text{woman}} \;\approx\; \mathbf{u}_{\text{queen}}
\]
\[
   \mathbf{u}_{\text{Paris}} - \mathbf{u}_{\text{France}} + \mathbf{u}_{\text{Italy}} \;\approx\; \mathbf{u}_{\text{Rome}}
\]

\medskip
\item Words living on the same semantic axis differ by approximately the same vector --- gender, capital city, comparative tense, monetary stance.

\medskip
\item \textbf{Bias is encoded too} (Bolukbasi et al., 2016, \emph{NeurIPS}; Caliskan, Bryson \& Narayanan, 2017, \emph{Science}):
\begin{itemize}
    \item $\mathbf{u}_{\text{man}} - \mathbf{u}_{\text{woman}}$ has a large component on \emph{programmer} vs.\ \emph{homemaker}.
    \item Word embeddings reproduce IAT-style biases at human magnitudes.
    \item For economic text: occupation/gender, race/credit-risk, region/sentiment correlations all transfer from the training corpus into the geometry.
\end{itemize}

\medskip
\item \textbf{Lesson:} embeddings are a mirror of the corpus. Validating that a sentiment vector reflects \emph{policy} stance and not \emph{speaker identity} is a real research step (returns in Topic 7.4).
\end{itemize}
\end{frame}

\begin{frame}{Limits of Static Embeddings}
\begin{columns}[T]
\column{0.48\textwidth}
\begin{tcolorbox}[colback=red!5, colframe=red!50!black, title=\textbf{What static embeddings miss}]
\begin{itemize}
\item \textbf{One word, many meanings (polysemy):} ``rate'' = interest rate? exchange rate? growth rate? One vector for all senses.
\item \textbf{Word order ignored:} ``GDP growth NOT slowing'' and ``GDP growth slowing'' have nearly identical mean vectors.
\item \textbf{Long-range context:} fixed window misses dependencies across paragraphs.
\item \textbf{Meaning drifts over time:} ``mask'' before vs.\ after 2020 are very different concepts.
\end{itemize}
\end{tcolorbox}

\column{0.48\textwidth}
\begin{tcolorbox}[colback=green!5, colframe=green!40!black, title=\textbf{What we need next}]
\begin{itemize}
\item Vectors that change with \emph{context}, not just word identity
\item A mechanism for tokens to ``look at'' other tokens, near or far
\item A way to learn arbitrary relational patterns, not just co-occurrence within a window
\end{itemize}
\medskip
$\Longrightarrow$ \textbf{The Transformer.}
\end{tcolorbox}
\end{columns}
\end{frame}

%=============================================================================
\section{Macro Applications}

\begin{frame}{From Embeddings to Macro Variables}
\begin{itemize}
\item Once a sentence (or document) is a vector, every downstream task is a familiar one:
\begin{itemize}
    \item \textbf{Classification:} hawkish / dovish / neutral $\Rightarrow$ logistic regression on the embedding.
    \item \textbf{Regression:} project the embedding onto a numeric outcome (sentiment score, uncertainty index).
    \item \textbf{Clustering / topics:} $k$-means or LDA on embeddings $\Rightarrow$ discover policy themes.
    \item \textbf{Retrieval:} nearest neighbor under cosine $\Rightarrow$ ``find me speeches like this one'' (Lecture 8).
\end{itemize}

\medskip
\item \textbf{Three concrete patterns from the literature:}
\begin{enumerate}
    \item \textbf{Signal} --- FOMC sentiment series predicting yields (Hansen \& McMahon, 2016, \emph{JIE}).
    \item \textbf{Measurement} --- ECB Word2Prices and central-bank communication tone (Acosta \& Meade, 2015; Schnabel et al., 2024).
    \item \textbf{Generation} --- LLM-summarized 10-K risk sections feeding a credit-spread regression.
\end{enumerate}

\medskip
\item \textbf{Topic 7.4} goes deep on each --- corpora (FOMC, SEC EDGAR, BIS speeches), domain models (FinBERT, CentralBankRoBERTa), validation, and pitfalls.
\end{itemize}
\end{frame}

\begin{frame}{Macro Application Gallery (Quick Tour)}
\begin{center}
\renewcommand{\arraystretch}{1.15}
\scriptsize
\begin{tabular}{|p{3.6cm}|p{4.4cm}|p{5.8cm}|}
\hline
\textbf{Application} & \textbf{Representation} & \textbf{Outcome / use} \\
\hline
EPU index (Baker--Bloom--Davis, 2016, \emph{QJE}) & Keyword counts & Monthly index, VAR shock \\
\hline
FOMC transparency (Hansen, McMahon \& Prat, 2018, \emph{QJE}) & LDA topics over transcripts & How disclosure regime changed speech \\
\hline
News-based recession (Bybee, Kelly, Manela \& Xiu, 2024, \emph{J. Finance}) & Topic shares from news & Recession nowcast vs.\ surveys \\
\hline
FOMC sentiment (Hansen \& McMahon, 2016) & Dictionary + Word2Vec & Predict Treasury yields, macro surprises \\
\hline
ECB Word2Prices & Static + contextual emb. & Map speeches to inflation expectations \\
\hline
FinBERT / CentralBankRoBERTa & Domain-tuned BERT contextual & Sentence-level stance on filings, minutes \\
\hline
\end{tabular}
\end{center}

\smallskip
\textbf{Common pattern:} (1) tokenize $\to$ (2) embed $\to$ (3) aggregate to document/series $\to$ (4) regress / forecast / classify. Steps (1)--(2) are \emph{this topic and its predecessor}; (3) is a research design choice; (4) is standard econometrics.
\end{frame}

\begin{frame}{Wrap-up: From Tokens to Transformers}
\textbf{What we covered.}
\begin{enumerate}\setlength{\itemsep}{1pt}
\item \textbf{Three roles of text:} signal, measurement, generation (Topic 7.1a).
\item \textbf{Tokenization:} BPE, WordPiece, SentencePiece. Subwords beat words and characters; tokenizer ships with the model (Topic 7.1a).
\item \textbf{From counts to vectors:} TF-IDF $\to$ Word2Vec / GloVe (static) $\to$ contextual embeddings.
\item \textbf{Geometry:} cosine similarity, analogies, bias --- all read off the embedding space.
\item \textbf{Macro:} every downstream task starts from a sentence vector; Topic 7.4 deep-dives.
\end{enumerate}

\medskip
\textbf{What's next: Topic 7.2a --- Attention Mechanism, then 7.2b --- Transformer Architecture.}
\begin{itemize}\setlength{\itemsep}{1pt}
\item Static embeddings can't see context, can't handle polysemy, can't combine across long ranges.
\item The Transformer fixes all three with one mechanism: \emph{self-attention} --- a learned, content-addressable lookup over the entire sequence.
\item Then: T3a unpacks training (pre-train $\to$ fine-tune $\to$ alignment); T3b audits limits; T4 surveys economic applications.
\end{itemize}

\medskip
\centering\emph{Tokens become vectors. Vectors become geometry. Geometry becomes economics.}
\end{frame}

\end{document}
